% 第3章3.4.3两阶段初始化（phase1b.py 验算，含词典序扰动）。三点算例加人工根 r，人工弧容量取总供给3、第一阶段费用1，原弧第一阶段费用0。
% (a) 星形初始树：s->r 流3，r->t 流3，a->r 流0；人工流总量6。
% (b) 入弧 s->t（约化费用-2），s->r 与 r->t 同时可降到0，词典序选 s->r 出基；再入 a->t（约化费用-2），零步长，a->r 出基。
%     第一阶段终止树 {s->t, r->t, a->t}，原弧流 (0,0,3)，人工流为0。
% (c) 删去 r 及人工弧，得原网络生成树 {s->t, a->t}，s->a 在 L。
\begin{tikzpicture}[>=Stealth,line width=.8pt,font=\fontsize{8}{10}\selectfont,
  v/.style={circle,draw=FigureStroke,fill=FigureBlueFill,line width=.8pt,minimum size=6.5mm,inner sep=0pt,text=FigureInk},
  rv/.style={v,fill=FigureYellowFill},
  tr/.style={->,draw=FigureBlue,line width=1.4pt,shorten >=.5pt},
  lo/.style={->,draw=FigureMuted,line width=.8pt,dashed,shorten >=.5pt},
  up/.style={->,draw=FigureMuted,line width=.9pt,densely dotted,shorten >=.5pt},
  lab/.style={text=FigureInk,inner sep=1.5pt},
  pot/.style={text=FigureGreen},
  note/.style={text=FigureMuted,font=\FigureNoteFont\fontsize{7.5}{9}\selectfont}]
  \foreach \dx/\k in {0/A,4.6/B,9.2/C}{
  \begin{scope}[xshift=\dx cm]
  \node[v] (s\k) at (0,0) {$s$};
  \node[v] (a\k) at (1.6,1.35) {$a$};
  \node[v] (t\k) at (3.2,0) {$t$};
  \end{scope}}
  % (a)
  \node[rv] (r1) at (1.6,-1.35) {$r$};
  \draw[lo] (sA)--(aA); \draw[lo] (aA)--(tA); \draw[lo] (sA)--(tA);
  \draw[tr] (sA)--node[lab,below left=-1pt]{$3$} (r1);
  \draw[tr] (r1)--node[lab,below right=-1pt]{$3$} (tA);
  \draw[tr] (aA) to[bend right=38] node[lab,left,pos=.3]{$0$} (r1);
  \node at (1.6,-2.3) {(a) 星形初始树，人工流$6$};
  % (b)
  \node[rv] (r2) at (6.2,-1.35) {$r$};
  \draw[lo] (sB)--(aB);
  \draw[tr] (aB)--node[lab,above right=-1pt]{$0$} (tB);
  \draw[tr] (sB)--node[lab,above]{$3$} (tB);
  \draw[lo] (sB)--(r2);
  \draw[tr] (r2)--node[lab,below right=-1pt]{$0$} (tB);
  \draw[lo] (aB) to[bend right=38] (r2);
  \node at (6.2,-2.3) {(b) 第一阶段结束，人工流$0$};
  % (c)
  \draw[lo] (sC)--node[lab,above left=-1pt]{$0$} (aC);
  \draw[tr] (aC)--node[lab,above right=-1pt]{$0$} (tC);
  \draw[tr] (sC)--node[lab,below]{$3$} (tC);
  \node at (10.8,-2.3) {(c) 删去人工根后的可行树基};
  \node[note] at (6.2,-2.9) {粗蓝线：树弧（旁注流量）；虚线：取下界的非树弧；黄色顶点为人工根};
\end{tikzpicture}
